\RequirePackage[T1]{fontenc}
\documentclass[10pt,a4paper]{article}
\usepackage[margin=20mm]{geometry}
\usepackage{amsmath,amssymb}
\usepackage{enumitem}
\usepackage[hidelinks]{hyperref}
\usepackage{mathptmx}
\setlength{\parindent}{0pt}
\setlength{\parskip}{4pt}
\setlist{nosep,leftmargin=14pt}
\newcommand{\loc}[1]{\textit{Manuscript: #1.}}
\newcommand{\GI}{G\textsubscript{1}}
\newcommand{\GII}{G\textsubscript{2}}
\newcommand{\GIII}{G\textsubscript{3}}
\newcommand{\GIV}{G\textsubscript{4}}
\newcommand{\comment}[2]{\par\medskip\noindent\textbf{#1}\hspace{0.4em}\textit{#2}\par}

\begin{document}
\begin{center}
{\large\bfseries Reviewer compliance report}\\[2pt]
WiCOMM-2026, Paper ID 167\\
\emph{Fail-Closed or Fail-Degraded? A Simulation Study of Forced Assurance Loss in
Contested Multi-Hop QKD Networks}\\[2pt]
Lakshya Jain, Siddhi Arora
\end{center}

We thank the reviewer for a careful and constructive review. The camera-ready
manuscript keeps the submitted title, research question, structure and core metrics,
adding floor-compliant service $Q$ and below-floor service $L$. It responds to every
main concern, minor point and question below; the one suggestion we did not follow in
full (point 9) is explained. Page numbers refer to the six-page camera-ready PDF. The
manuscript itself presents the corrected results without referring to the review;
this report records what changed and why.

\section*{Corrections found during revision}
While re-running the study we found and corrected the following. Each is backed by
runs that are released with the paper.

\begin{itemize}
\item \textbf{Stored key after abort.} The submitted simulator treated key already
stored on a link as unusable once the link aborted generation. Stored key does not
expire when generation stops. With every other submitted setting and identical random
streams, correcting this one rule moves fail-closed availability at $k=2$ from 0.304
to 0.767, and forced assurance loss (as defined in the submission, over served
entitled requests) from 67.4\% to 4.3\%; reverting it reproduces the submitted values
to within 0.07 points. The submitted cliff therefore measured the loss of stock, not
of key generation. The camera-ready model keeps stored key spendable, declares the
stock through an initial usable fill (2\% in the main runs), and reports stock and
attack duration as part of the forcing cost (\S V-D).
\item \textbf{The submitted $k=2$ attack was not the minimum cut.} The static-betweenness adversary bought BDE1--BDE2 and BDE1--REL1, which isolate two of the
three forward posts. Table II now reports this adversary beside the cost-weighted
minimum cut; with stock retained, its first link costs 0.7 points of fail-closed
service, so the submitted rise at $k=1$ disappears.
\item \textbf{Formation size.} The formation has eight nodes and nine links, not
nine nodes.
\item \textbf{Chain forcing cost.} On the chain every link lies on one cycle with
equal betweenness, so the static ranking is decided by tie order; with our tie
order it needs $d+2$ links (the submitted order needed $d+1$). Residual betweenness
needs two at every depth $d\ge2$ (\S V-C).
\item \textbf{Key-aware routing.} The submitted 2.7-point fail-closed loss of
key-aware routing reproduces, but only in a window after a 12\,h warm-up from full
stores; from $t=0$ or from 2\% fill key-aware routing gains 0.08--0.13 points
(\S V-B).
\item \textbf{Withdrawn claims.} Class-aware reservation floors are tunable
(fail-closed service 96.6, 81.2 and 59.9\% at floor scales 0.05, 0.25 and 1), and the
per-class shield is no longer circumvented through omitted classes; both submitted
statements are withdrawn. The rate formula's ITS zero at $f=1.16$ is 9.8\%, not
9.9\%, so at $e_{\mathrm{abort}}=0.11$ complete suppression needs $\Delta e=0.068$,
not 0.08.
\end{itemize}

All results are regenerated with the corrected code; 127 regression tests and the
43-check validation gate pass.

\textbf{Settings changed at the reviewer's request} (point numbers refer to the
review): $n_{\mathrm{auth}}=7$ in all results, with 2 and 40 as sensitivities
(point 4); ISR arrival rate rescaled from 0.02 to 0.002\,s$^{-1}$, so the mean debit
across the formation's cut is 83\% of its supply, with the submitted rate kept as an
overload case (point 8); per-link attack costs (point 5); weather-matched baselines
(point 6). Main comparisons use twelve seeds, sensitivities eight, and the online
adversary and abort-threshold sweep six, with 95\% percentile bootstrap intervals
over seeds (20\,000 resamples) and common random numbers for paired differences.

\section*{Main concerns}

\comment{1. Single topology.}{All headline numbers come from one nine-node graph \ldots\ Evaluate on multiple topologies with varying cut size and report how the metric scales.}
We added four graphs---an eight-node ring, a $3\times3$ grid and three- and
four-path parallel graphs---giving minimum cuts of two, two, three and four links
beside the formation (\S III, Fig.~2). We agree that the headline numbers are a
property of one graph; the paper now says so and reports how the result scales.
With the budget equal to the cut, forced assurance loss is 87.6--92.9\% on all five
graphs, and fail-closed service follows a stock-limited prediction from a new cut
condition, Eq.~(4): 6.0/9.0/12.0\% predicted for cuts of 2/3/4 against
6.5--6.7/9.6/12.3\% observed. With one cut link suppressed, loss appears only where
the remaining cut links cannot carry the demand (4.8--11.5\% on two-link cuts, 0.71\%
and 0.05\% on three- and four-link cuts). The conclusion states the five-graph result.
\loc{\S III and \S III-C, Eq.~(4), pp.~2--3; \S V-A and Fig.~2, pp.~3--4; \S VII, p.~6}

\comment{2. Non-equivalent metrics.}{Under fail-degraded, a session counts as available even when served with post-quantum crypto only \ldots\ make the asymmetry explicit in the abstract.}
Of the two options offered, we took the second, and added a common denominator: timely service $A$, floor-compliant
service $Q$ and below-floor service $L$ share one arrival denominator, $A=Q+L$
(Eq.~3), and every headline comparison reports $A$ and $Q$ together. The abstract
states the asymmetry directly: fail-degraded operation keeps 97.5\% of requests
served but only 6.6\% at their authorised grade, against fail-closed service of
6.7\%. Forced assurance loss keeps blocked entitled requests in its denominator, so
both modes are scored against the same entitled arrivals. The composite
$J=A-0.75L$ is used only in the policy audit, where the value of below-floor service
is an explicit choice.
\loc{Abstract; \S III-C, Eq.~(3), p.~3; \S V-A and Table II, pp.~3--4}

\comment{3. Validation overstated.}{V1 compares the model against its own accounting; V4 bounds only to an order of magnitude \ldots\ State what each check establishes.}
We withdrew the claim of two reproductions. V1 is described as arithmetic
consistency, V4 as a bound to an order of magnitude, and V2 and V3 as two failures in
the same direction, which fix the direction of the accounting bias and motivate the
new default $n_{\mathrm{auth}}=7$. The abstract no longer mentions validation.
\loc{\S IV-A, p.~3}

\comment{4. Wrong default parameter.}{Field data implies $n_{\mathrm{auth}}$ of 7 to 40 \ldots\ Report main tables at the calibrated value with 2 as the optimistic sensitivity.}
All results now use $n_{\mathrm{auth}}=7$ (a \GII\ debit of 1385\,bit per hop), with 2
and 40 as sensitivities; the one exception, labelled in the text, is the stored-key
contrast run at the submitted settings. As the reviewer anticipated, the choice is
load-bearing where scarcity binds: under the cut, fail-closed service is
8.5/7.0/3.3\% and forced loss 87.9/89.9/93.5\% at $n_{\mathrm{auth}}=2/7/40$ (eight
seeds). We withdrew the claim that policy differences cancel as a common additive
term. The \GI\ affordability and routing results were re-run at the calibrated value.
\loc{\S III-B, p.~3; \S V-D, p.~5}

\comment{5. Uniform attack cost.}{The topology mixes fibre and free-space links, yet budget $k$ assumes all links are equally attackable.}
Each link now carries a relative attack cost, 1 for free-space and 10 for the fibre,
with ratios 1 and 100 as sensitivities, and budgets are spent in these units. The
paper states that the costs are scenario weights, not measured expenditure. The
results are cost-conditioned: with the fibre 10 or 100 times dearer, a
residual-betweenness adversary avoids it and forces 10.1\% at budget two, against
89.8\% when all links cost the same, while the cost-weighted minimum cut, which
contains no fibre, is unaffected (eight-seed sensitivity; ratios 10 and 100 give
identical results).
\loc{\S III-A, p.~2; \S V-C, p.~4}

\comment{6. Correlated outage affects the baseline.}{Weather or terrain masking could sever the min-cut with no adversary present \ldots\ Address the implication for the zero baseline.}
We agree, and no longer claim that all loss under attack is attack-attributable.
Attack effects are reported as paired differences between attacked and unattacked runs under
identical environmental draws. With up/down times of 900/600\,s the unattacked forced
loss is 1.7\% (independent links) and 3.0\% (all free-space generation failing
together), while the paired attack excess stays at 88.1 and 86.9 points. The paper
states that $F$ counts loss whatever its cause, that attribution needs the
weather-matched baseline, and that ``a held correlated outage looks like an attack''.
\loc{\S IV, p.~3; \S V-D, p.~5}

\comment{7. Metric ignores exposure.}{A downgraded 2\,kbit command message is weighted the same as a 320\,kbit logistics transfer \ldots\ Add weighting by data sensitivity.}
We added a payload-weighted forced loss $F_B$ and class-importance weights $(5,2,1)$
and $(10,3,1)$ for (C2, ISR, LOG), labelled as exposure proxies rather than measured
harm. At the headline cut the payload-weighted forced loss is 90.3\% against 90.1\%
by count. In the policy audit the weighting matters: the learned controller's small
utility gain comes with $+1.7$ points of below-ITS payload exposure, which the
count-based $\Delta F$ ($+0.33$ points) understates; class weighting leaves its gain
unchanged.
\loc{\S III-C, p.~3; \S V-A; \S V-E and Table IV, p.~5}

\comment{8. Rescale the \GI\ class.}{A class that cannot be served under any configuration distorts the grade mix and the routing comparison. Rescale ISR to what the links can carry.}
ISR is rescaled to 0.002\,s$^{-1}$, so the mean debit across the formation's cut is
397\,bit/s, 83\% of its 480\,bit/s supply; the submitted rate (over six times the
supply) is kept as an overload case and 0.03 of it as light load. At the rescaled load
ISR reaches its one-time-pad floor in 97.0\% of arrivals, against 12.6\% at the
submitted rate, so the submitted verdict that the floor was unimplementable is
corrected: a \GI\ floor fails through overload, not by construction. The routing
comparison is re-run at the rescaled load.
\loc{\S III-B, p.~3; \S V-A, p.~4; \S V-B, p.~4}

\comment{9. Policy search is not required.}{The useful output is the threshold, which follows from the objective definition \ldots\ ``104.8\%'' bounds only the search budget. Compress \S V-E.}
We agree that the crossing price is algebra on measured outcome differences, and the
paper now presents it that way, reporting the $\lambda$ threshold directly: in
$J=A-\lambda L$ the heuristic beats the strict floor rule for $\lambda<0.98$ (CI
$[0.96,1.00]$). The 104.8\% claim is removed, and the paper states that the
heuristic's share of the learned policy (now 99.8\%) reflects the search budget, not
the headroom. We compressed \S V-E rather than removing it: its text is shorter than
in the submission (about 235 against 250 words), its numbers moved into the new
Table IV, and with that table it occupies the same space as before; the room for the
topology sweep came from compressing \S I--\S IV.

We kept a compact learning study because it now supports a claim the objective alone
does not give, and one the paper's main result depends on. Where the reviewer's
$\lambda$ algebra ranks fixed policies, a learned controller shows what an optimiser
does with the shield's exception allowance: it spends it where it costs the most
one-time-pad payload (below-ITS payload exposure $+1.7$ points, ISR compliant service
$-2.8$), a gain invisible to $\lambda$ because $\Delta L$ is indistinguishable from
zero and reversed only by pricing forced loss ($\mu^\ast=0.32$). This is the policy
counterpart of the paper's argument that assurance loss must be measured, not
inferred from availability. It also shows that tables trained on nominal traffic
alone fail on stressed conditions ($-4.7$ and $-3.6\times10^{-3}$), because nominal
runs never reach depletion. The protocol now uses disjoint training, validation and
test seeds, frozen tables, a per-class shield and seven test conditions including two
unseen graphs. If the reviewer prefers, \S V-E can be reduced to its first paragraph
and Table IV without affecting any other result.
\loc{\S V-E and Table IV, p.~5}

\section*{Minor points}

\comment{$\gamma=3.4$ is fitted to one datapoint from a different QKD modality.}{Call it an illustrative bound, not a calibration.}
$\gamma=3.4$ is now described as an illustrative steeper response, not a calibration.
The new sub-threshold results show where it matters: below the abort threshold a steep
response can force loss without aborting any link (Table III).
\loc{\S III-A, p.~2; Table III, p.~5}

\comment{\S V-B(1): the zero-width CI follows from the setup.}{Soften the wording.}
The paper now states that the zero-width interval is an identity of the model, not a
measured precision, and that floor-compliant service $Q$ still discriminates between
the routing policies.
\loc{\S V-B(1), p.~4}

\comment{Table III needs confidence intervals; Table II needs a per-class breakdown and how SD $\le0.012$ was computed.}{}
The former Table III is replaced by a capability-against-budget table over three
graphs, which reports every non-zero cell with its 95\% bootstrap interval. Table II
gives floor-compliant service per class (C2, ISR, LOG), and its caption states the
largest bootstrap half-width and the largest sample standard deviation across seeds,
computed per cell over the twelve seeds. The blanket ``SD $\le 0.012$'' statement was
removed.
\loc{Tables II and III, pp.~3 and 5}

\comment{Table I: the all-\checkmark\ row for ``This work'' is self-serving.}{State the gap in prose.}
The row is removed; the gap is stated in the Related Work prose, scoped to the studies
we located. We also added prior work found during revision: denial-of-service key
depletion (Mehic \emph{et al.}, 2022) and single-link studies of attack- or
failure-triggered KEM fallback (Agalya and Priyadarsini, 2026; Hernandez-Hernandez
\emph{et al.}, 2026).
\loc{\S II and Table I, pp.~1--2}

\comment{Fig.~1(b) is overloaded.}{Separate the two attack surfaces.}
Fig.~1 is redrawn. Panel (a) shows the eight-node formation and the minimum cut that
isolates its forward posts; panel (b) lists the three failure surfaces separately---
quantum-channel attack, classical-transport denial and online key-establishment
denial---with what each stops and what still works.
\loc{Fig.~1, p.~2}

\comment{Nearly every clause carries a caveat.}{Compress and use the space for the missing topology sweep.}
Caveats are stated where a choice is introduced and collected once in the limitations
paragraph (\S VI); for example, attack costs and payload weights are declared as
scenario choices. The space is used for the topology sweep, the cut condition and the
answers to the questions below.
\loc{\S VI, p.~6}

\comment{Reference [5] is dated 2026.}{Confirm publication status.}
Confirmed: P. Spooren, A. Neuhold, S. Ramacher and T. H\"uhn, ``PQC-enhanced QKD
networks: A layered approach,'' \emph{Proc. IEEE QCNC}, Kobe, Japan, 2026,
pp.~354--358, doi:10.1109/QCNC69040.2026.00060. All other references were checked
against publisher records; several were updated to their published versions (for
example, the IETF hybrid key exchange draft is now RFC~9954).

\section*{Questions}

\comment{Q1. If ISR demand is scaled to what the free-space hops can sustain, does the fail-mode gap persist, narrow, or invert?}{}
It persists under attack, narrows modestly, and never inverts. Unattacked, the gap
vanishes at sustainable load (paired $\Delta A$ of 0.00 and $+0.01$ points at 0.03 and
0.1 of the submitted rate) and appears only through overload at the submitted rate
($+17.2$). Under the cut, $\Delta A$ is 85.4, 90.6 and 94.7 points at 0.03, 0.1 and 1
of the submitted rate (eight seeds; 90.8 on the twelve main seeds), narrowing only
because a lighter load drains the stock more slowly, and fallback never raises the
formation's compliant service. The size of the gap is set by the stock behind the
cut and the attack's duration: from full stores the same 24\,h attack forces nothing.
\loc{\S V-D, p.~5}

\comment{Q2. Does the capability/budget non-substitution hold for graphs with min-cut greater than 2?}{}
For the linear response ($\gamma=1$) and without outages, yes: on cuts of two, three
and four links, sub-threshold injection forces under 0.1\% at any budget, and a budget
one short of the cut forces under 0.1\% at any injection; only a full cut past the
threshold forces loss (93.1, 90.4 and 87.7\%), and every budget from the cut upward
gives the same value. With outages, a budget one short of a two-link cut forces
4.8--11.5\% (Fig.~2b), because the remaining link cannot carry the demand. For the
steep illustrative response ($\gamma=3.4$), no: sub-threshold injection lowers the
cut's supply toward what the entitled classes need and forces 63.0, 41.6 and 17.3\% at
$\Delta e=0.03$ without aborting any link. Table III reports every non-zero cell with
its interval. We also corrected the submitted statement about raising the abort
threshold to 0.11: complete suppression then needs $\Delta e=0.068$, not 0.08,
because the rate formula's ITS zero (9.8\%) binds before the abort threshold.
\loc{\S V-D and Table III, pp.~4--5}

\comment{Q3. Under an adversary that re-targets as the router reroutes, does the centrality-guided forcing cost still grow linearly with depth?}{}
Not for a competent adversary. The linear growth came from a static ranking that, on
the chain, is decided by tie-breaking, since every chain edge lies on one cycle with
equal betweenness. An adversary that recomputes betweenness on the residual graph needs
two links, the minimum cut, at every depth beyond one. We also tested an online
adversary that sees only committed routes and updates every 60\,s (depths 2, 4 and 6):
if it keeps its slots it holds a cut from its first choice and forces 89.7--92.9\%
with 48 link-hours; if it releases slots whenever its path is momentarily
unreachable, its forced loss falls with depth (72.3\% to 35.5\%), so depth protects
only against that less persistent variant.
\loc{\S V-C and Fig.~3, p.~4}

\comment{Q4. Would the classical-channel dependency change if the fallback used pre-placed symmetric keys instead of an online KEM?}{}
Not the transport dependency: denying the classical channel on the cut drops service
to 0.06\% with either fallback, because ciphertext still needs a path. Pre-placed keys
remove a different dependency, on a live key-establishment service: with that service
down under the held cut, an online KEM fallback serves 5.0\% and pre-placed keys
97.5\%. They still force 89.9\% of entitled requests below ITS, so they relocate the
dependency without reducing assurance loss. The model treats pre-placed keys as ideal;
their lifetime, refresh and distribution cost are stated as unmodelled.
\loc{\S III-A and Fig.~1(b), p.~2; \S V-D, p.~5}

\section*{Reproducibility}
The code, configurations, per-seed results, the scripts that generate every table and
figure, and the experiment that compares the submitted and corrected simulators on
identical random streams are in the project repository. The manuscript's numerical
tables are generated from raw run files by a single script, which also records each
cited value with its cases and seeds.
\end{document}
